% !TeX program = XeLaTeX
% !TeX root = ../../shloka.tex
% --meta--
% deity: Guru Parampara
% composer: Sadashiva Brahmendra
% type: Mala
% wiki_title: Jagadguru Ratna Mala
% --end-meta--

\sect{जगद्गुरुरत्नमाला}

\twolineshloka
{यद् अबोधवशाद् अहं ममेदं तद् इहेत्यादिरुदेति भूरिभेदः}
{तद् अखण्डम् अनन्तम् अद्वितीयं परम् आनन्दमयं पदं श्रयेयम्}% ॥ १ ॥

\twolineshloka
{वटपादपमूलपट्टभद्रं स्फुटचिन्मुद्रम् उदूढमौनमुद्रम्}
{कलितेन्दुकलं कृपासमुद्रं कलये कञ्चन देशिकं विनिद्रम्}% ॥ २ ॥

\twolineshloka
{कमलाकुचकुम्भकुङ्कुमश्रीकलितोरःस्थलम् अप्यरागयोगम्}
{करवै मुरवैरिणं महान्तं करुणार्द्रं हृदये कविं पुराणम्}% ॥ ३ ॥

\twolineshloka
{गरुडाङ्कगभीरनाभिपद्मोदरभूगह्वरनिष्ठनैष्ठिकाग्र्यम्}
{चतुराननम् आश्रयामि कामं चतुराम्नायवदावदं सदाऽपि}% ॥ ४ ॥

\twolineshloka
{कुशिकप्रभवास्त्रशस्त्रपालीप्रशमाचार्यकसिद्धहस्तदण्डम्}
{शमवन्तम् अरुन्धतीसहायं शरणं यामि वसिष्ठम् अप्रमेयम्}% ॥ ५ ॥

\twolineshloka
{बहुमित्रसहक्षितीशपापापहसान्निध्यम् अरोध्यसत् प्रतापम्}
{महद् अञ्चितमन्त्रयन्त्रशक्तिं मनसा शक्तिम् उपैमि सद् विरक्तिम्}% ॥ ६ ॥

\twolineshloka
{पवनाकुलकाशलेशकल्पं प्रभवादेव पराहताशराशम्}
{प्रणमामि पराशरं मुनीशं प्रचरद् वैष्णवसंहिताप्रकाशम्}% ॥ ७ ॥

\twolineshloka
{निगमान् अपि योऽन्वशाच्चतुर्धा व्यधिताष्टादशधाऽपि यः पुराणम्}
{स च सात्यवतेय ईप्सितं मे सकलाम्नायशिरोगुरुर्विधत्ताम्}% ॥ ८ ॥

\twolineshloka
{जननीजठराद् इव च्यवन् यो जगतो नाद्रवद् आत्मविद् विपद्भ्यः}
{अनहन्तम् अहं तम् आत्मवन्तं भगवन्तं शुकम् आश्रये प्रशान्तम्}% ॥ ९ ॥

\twolineshloka
{अभियुञ्जद् अयार्च्यपूज्यपादान् अपलून्यादिनिषाकसिद्धनेतॄन्}
{अथ गौडपदान् फणीशभाष्यप्रथमाचार्यकपण्डितान् प्रपद्ये}% ॥ १० ॥

\twolineshloka
{हरितल्पहराङ्घ्रिनूपुरक्ष्माधरसौमित्रिबलात्रिपुत्रलक्ष्मा}
{जयताद् उपरेवम् आत्तधामा जयगोविन्दमुनिः स चन्द्रनामा}% ॥ ११ ॥

\twolineshloka
{प्रभवात् परमार्थसङ्ग्रहाच्च प्रणयस्थानम् अमुष्य सिद्धशिष्यः}
{विरतक्षितिरक्षणो विरक्त्या हरिरंहो मम तृंहतां स्वशक्त्या}% ॥ १२ ॥

\twolineshloka
{कलिना बलिनाऽखिले खिलेऽपि स्खलिते श्रौतपथेऽपथे प्रवृद्धे}
{जपहोमतपस्सु नामशेषेष्वपि यातेषु सुभाषितेषु शोषम्}% ॥ १३ ॥

\twolineshloka
{जिनबुद्धकपालिशक्तिवह्नीन्द्विनसेवासु समन्ततः श्रितासु}
{शिथिले द्विजवेदवेदिवादे प्रथिते चाद्वयवेदबाह्यमोदे}% ॥ १४ ॥

\twolineshloka
{जगद् ईक्षणविह्वलामृतान्धोनिगदव्यक्तकृपारसानुबन्धम्}
{प्रणिदिश्य गुहं पुरैव गन्तुं प्रणिबन्धुं च मखान् द्विषश्च यन्तुम्}% ॥ १५ ॥

\twolineshloka
{अवतार्य सुरान् परांश्च पूर्वं विधिविष्ण्विन्द्रमुखान् विनोदपूर्वम्}
{स्वयम् अप्यवतीर्य सुत्युरार्याकमितुः श्रीशिवशर्मणो विचार्य}% ॥ १६ ॥

\twolineshloka
{उदभूत् सदने निटालदृग् यो मदभाजां सुधियां प्रमाथयोग्ये}
{शिशुरर्पयतान्मुमुक्षुभाग्यं स शुभं शङ्करदेशिकः सुभोग्यम्}% ॥ १७ ॥

\twolineshloka
{उपनीय दिवं गते स्वताते व्यपदिश्यागमसङ्ग्रहेऽप्यतीते}
{निजमातरम् अन्वमीमनद् यो नियमाप्तिं स मुदे ममास्तु सद्यः}% ॥ १८ ॥

\twolineshloka
{प्रतिचन्द्रभवं निवृत्तिधर्मा श्रितगोविन्दमुनेरवाप्तधर्मा}
{जयतात् कृतसूत्रभाष्यकर्मा स्वयम् अन्तेवसतां वितीर्णशर्मा}% ॥ १९ ॥

\twolineshloka
{कुहनान्त्यजविश्वनाथसृष्टो द्रुहिणव्यासवरोदितानुशिष्टः}
{ममतां मम तावद् एष भिन्द्यान्नमतश्चोपरतिं ददात्वनिन्द्याम्}% ॥ २० ॥

\twolineshloka
{अपि हालनृपालपालितं द्राग् अपहायेन्द्रहरिन्मुखं वितन्द्रम्}
{श्रितमुर्मुरम् उद्धरंस्त्रिवेण्यां यतिराड् भट्टम् अघानि वो वितृण्यात्}% ॥ २१ ॥

\twolineshloka
{पटुमण्डनमिश्रखण्डनार्थं प्रविशन् पद्मवनं नवं जयार्थी}
{तम् अधृष्यगिरं विधृष्य वादैर्यतिम् आधात् स च पातु मां प्रमोदी}% ॥ २२ ॥

\twolineshloka
{प्रविशन् बदरीम् अवाप्य सद्यः परमाचार्यपदार्चनं क्रमाद् यः}
{धवलाचलम् आप्य योऽप्यमाद्यच्छिवलावण्यम् उदीक्ष्य तं प्रपद्ये}% ॥ २३ ॥

\twolineshloka
{प्रतिपादितलिङ्गपञ्चकेऽमुं प्रणिवर्त्याशु तिरोहिते गिरीशे}
{विनिवृत्य स दिग् जयप्रवृत्तो विविधैः शिष्यवरैर्विभातु चित्ते}% ॥ २४ ॥

\twolineshloka
{कलयन् निलयं च तुङ्गभद्रातटिनीरोधसि वेधसः स्त्रिया द्राक्}
{कतिचिच्छरदोऽत्यवीवहद् यो यतिराट् क्वापि मठे स मेऽस्तु सद्यः}% ॥ २५ ॥

\twolineshloka
{अथ कान्यकुमारसन्धिसेतुस्थलिनीवैङ्कटकालहस्तियातुः}
{यमिनेतुरमुष्य काञ्चियात्रा शम् इदानीं शमदं क्रियाद् विचित्रा}% ॥ २६ ॥

\twolineshloka
{श्रितनिर्मलराजसेनचोलक्षितिपालोद्धृतविप्रदेवशालः}
{वरदस्य तथाऽऽम्रनायकस्याप्युरुवेश्मद्वयकृज्जयाय मे स्यात्}% ॥ २७ ॥

\twolineshloka
{प्रकृतिं च गुहाश्रयां महोग्रां स्वकृते चक्रवरे प्रवेश्य योऽग्रे}
{अकृताश्रितसौम्यमूर्तिम् आर्यां सुकृतं नः स चिनोतु शङ्करार्यः}% ॥ २८ ॥

\twolineshloka
{परकीयवपुःप्रवेशशैल्या स्मरतन्त्राणि विदन् वशी न लौल्यात्}
{अकृतातनुभारतीं वशे यः सुकृतं मे सुतरां गुरुः स चीयात्}% ॥ २९ ॥

\twolineshloka
{उपयात्सु बुधेषु सर्वदिग्भ्यः प्रदिशन्नाशु पराभवं य एभ्यः}
{विधृताखिलवित्पदश्च काञ्च्याम् अधृतार्तिः स दिशेच्छ्रियं च काञ्चित्}% ॥ ३० ॥

\twolineshloka
{समतिष्ठिपद् आहिमाद्रिसेव्यं क्रमशो धर्मविचारणाय दिव्यम्}
{अधिकाञ्चि च शारदामठं योऽभ्यधिकं नः सुखम् आतनोतु सोऽयम्}% ॥ ३१ ॥

\twolineshloka
{पृथिवीधवपद्मपादहस्तामलकानन्दगिरीन् अतिप्रशस्तान्}
{चतुरोऽकृत यश्चतुर्षु पीठेष्वतुलां नश्चतुरो गिरं स देयात्}% ॥ ३२ ॥

\twolineshloka
{परम् अन्तिकसत् सुरेश्वराद्यैः परमाद्वैतमतं स्फुटं प्रवेद्य}
{परिकाञ्चिपुरं परे विलीनः परमायास्तु शिवाय सद्गुरुर्नः}% ॥ ३३ ॥

\twolineshloka
{स्थिरबोधघनप्रतापदाम्नोर्गुरुपृथ्वीधवविश्वरूपनाम्नोः}
{चिरम् अर्थनयोपतुङ्गभद्रं सरसः सौतु सुरेश्वरः स भद्रम्}% ॥ ३४ ॥

\twolineshloka
{प्रणयन् लययोगम् अक्षुद् अञ्च्यं पुनराचार्यमठे निपद्य काञ्च्याम्}
{भुवि यः शिवलिङ्गभूयमन्ते प्रविवान् सोऽस्तु सुखाय नः सुशान्तः}% ॥ ३५ ॥

\twolineshloka
{प्रथिते किल कामकोष्ठनाम्नि प्रथते पुण्यरसेति यत् पवित्रम्}
{पदमस्य गुरोः समाधिभूमेः प्रणये मण्डनचिह्नितं निकामम्}% ॥ ३६ ॥

\twolineshloka
{जपद् इष्टदकामकोटिपीठाधिपभूमण्डलदेशिकेन्द्रपाठात्}
{नरनुत्यशुचीनि निर्णिनिक्षे परहंसस्तवतोऽक्षराण्यरूक्षम्}% ॥ ३७ ॥

\twolineshloka
{सलिलाशन एव यः सलीलं विलयं प्रापिपद् आर्हतान् सुशीलः}
{सुमहारफणीन्द्रयोः सदृष्टिः स हि सर्वज्ञगुरुर्ह्रियात् कुदृष्टिम्}% ॥ ३८ ॥

\twolineshloka
{यमिपञ्चकम् अञ्चितोज्ज्वलाख्यं शमितान्योक्ति च सत्यबोधमुख्यम्}
{परवैभववारिविद्यम् अन्तः करवै तत् त्रितयं च चिद्घनान्तम्}% ॥ ३९ ॥

\twolineshloka
{सुमहाव्रतसार्वभौमचर्याश्रमहानाय निगूढयोगचर्यः}
{शमदां मम चन्द्रशेखरेन्द्रः स मुदं पातु जगद्गुरुर्वितन्द्रः}% ॥ ४० ॥

\twolineshloka
{फलहायनकाष्ठमौनसेवाबलहासास्पदबालमुग्धभावम्}
{अजरामरम् आत्तलिङ्गभावं भज सच्चिद्घनम् अन्तरङ्ग देवम्}% ॥ ४१ ॥

\twolineshloka
{कलया मलयाद्र्यगस्त्यकूटस्थलयातूयितभैरवप्रमोटम्}
{कलयाम लयाध्वलीढभानुं चलविद्याघनम् आप्तकामधेनुम्}% ॥ ४२ ॥

\twolineshloka
{गुरुमार्गगवेषणात्तवेलाचलदृष्टाग्रजशिष्टमन्त्रलोलम्}
{हरपूजनपूतम् अस्तमोहं वरगङ्गाधरगीष्पतिं श्रयेऽहम्}% ॥ ४३ ॥

\twolineshloka
{कुलशेखरसाहितीप्रणेत्रे वलद् आसेतुहिमाद्रिविज्ञनेत्रे}
{जरदृष्टिशकादिघृष्टिदात्रे स्थिरवर्णाश्रमधर्मसंविधात्रे}% ॥ ४४ ॥

\twolineshloka
{स्तनितं निगमेषु शङ्करेण ध्वनिम् अद्वैतम् उदस्यते चिरेण}
{प्रणयानि महायतीशभूम्ने प्रणतीरुज्ज्वलशङ्करार्यनाम्ने}% ॥ ४५ ॥

\twolineshloka
{अथ गौडसदाशिवेन्द्रपादान् प्रथितान् न्यक्कृतबाह्लिकप्रवादान्}
{जननक्षण एव जातबोधान् जगद् आचार्यपदान् श्रये विबाधान्}% ॥ ४६ ॥

\twolineshloka
{विवदन् धिषणेन साधु दुर्दीदिविसाह्याय समागतेन सार्धम्}
{जयिजीर्णसुरेन्द्रभूमिभोगी जयतान्नः सुधिये सुरेन्द्रयोगी}% ॥ ४७ ॥

\twolineshloka
{चिरसाधितषोडशीप्रसादक्षरितश्वित्ररुजं निरस्तवादम्}
{स्ववशासुविसर्गम् अर्कभृत्यं नवविद्याघनम् आश्रितोऽस्मि नित्यम्}% ॥ ४८ ॥

\twolineshloka
{घनवाग् जनिमूक एव विद्याघनकारुण्यत ऊढहृद्यविद्यः}
{कलयन्नपि मेण्ठकं कवीन्द्रं कलयेन्नः स शिवानि शङ्करेन्द्रः}% ॥ ४९ ॥

\twolineshloka
{अपि यः श्रितमातृगुप्तविद्याधिपसेतुप्रवरादिसूरिहृद्याम्}
{सुषमाम् अधिताहिमाद्रिभूमौ विषमादित्यनुतोऽवतात् स चामुम्}% ॥ ५० ॥

\twolineshloka
{अनुजह्नुसुतातिवाहिताहं विधिनाऽऽराधितशङ्करं विमोहम्}
{शरणं समुपैमि चन्द्रचूडं चरणप्रौढतपःपतत्रिनीडम्}% ॥ ५१ ॥

\twolineshloka
{परिपूर्णविबोधनामभाजां धुरि गोकर्णधराभुवां विराजत्}
{मम मानसम् अस्यसिद्धिभूम्ना ममतार्तिं विजहातु मोदसीम्ना}% ॥ ५२ ॥

\twolineshloka
{अपि नास्तिकम् आस्तिकं दधानान् लपितैर्ज्यौतिषम् आर्यभट्टम् ऊनम्}
{प्रतिषिद्धपयोधिपोतयानान् यतिसच्चित्सुखसद्गुरून् नमामि}% ॥ ५३ ॥

\twolineshloka
{समधिष्ठितकोङ्कणे चिरेण स्वधृते साधु समाधिनाऽदरेण}
{निरतो नियमेन चित्सुखेऽस्मि श्रितसच्चित्सुखलक्षणे परस्मिन्}% ॥ ५४ ॥

\twolineshloka
{कृमिकीटपतङ्गसत्त्वभाषाक्रमविज्ञः स धुनोतु सत्त्वदोषान्}
{रससाधनसिद्धलिङ्गरूपः स चिदानन्दघनः ससत् स्वरूपः}% ॥ ५५ ॥

\twolineshloka
{यमिधीघनचिद्विलाससाधू समहादेवसमग्रबोधबोधौ}
{भवभूतिकृतार्चनौ च वन्दे सुखसद् ब्रह्मचिद् उज्ज्वलौ घनेन्द्रौ}% ॥ ५६ ॥

\twolineshloka
{स्वमठस्थम् अधीतभूरिभाषं क्रमिकार्चाहितकामदृक् प्रतोषम्}
{प्रचितस्वमठालयं च सच्चिद् रुचिरानन्दघनेन्द्रम् आश्रयेऽहम्}% ॥ ५७ ॥

\twolineshloka
{अपरानपि चन्द्रशेखरेन्द्रान् विपुलक्ष्माभ्रमणे निरस्ततन्द्रान्}
{शरणागतसर्वदुःखहर्तॄन् शरणं यामि कुवादिजैनजेतॄन्}% ॥ ५८ ॥

\twolineshloka
{भज रे बहुरूपचित्सुखेन्द्रं बहुशः सह्यगुहातपस्स्वतन्द्रम्}
{प्रशमप्रणयात् प्रशान्तचारं परम् आनन्दयुजं च पापदूरम्}% ॥ ५९ ॥

\twolineshloka
{प्रचिते परितस्तुरुष्कचक्रे निचिते म्लेच्छगवीविभूम्नि वक्रे}
{कथमप्यवताऽऽश्रितार्यवृत्तान्यथ विद्याघनमौनिनैमि तृप्तिम्}% ॥ ६० ॥

\twolineshloka
{निजतातमृतेः परं प्रकाशे गजगर्भे त्रिशरद् विपच्यकोशे}
{गमयन्तम् इवोल्बवासशेषं गणये विश्वजितः सुतं विशेषात्}% ॥ ६१ ॥

\twolineshloka
{परिवादभयात् प्रसूविसृष्टं विपिने व्याघ्रवधूभिरेव पुष्टम्}
{नवशङ्करम् आदिमाच्च धृष्टं नम माध्यन्दिनिनैव बाल्यशिष्टम्}% ॥ ६२ ॥

\twolineshloka
{परम् अभ्रतनोर्गिरैव शम्भोर्वरविद्याघनविष्टरोपलम्भी}
{अधिकाञ्चिपुरीमठं शुभं भोस्त्र्यधिकश्चित्त तनोतु ते विदम्भः}% ॥ ६३ ॥

\twolineshloka
{विदिशश्च दिशश्च सञ्चरन् यः प्रदिशन् भाष्यम् अदूष्यम् आस्त धन्यः}
{तम् अशेषविद् आसनाधिरूढं नम भट्टोद्भटवाक्पतेः सनीडम्}% ॥ ६४ ॥

\twolineshloka
{श्रुतियुक्तिशरैः परान् प्रणिघ्नन् इति चाद्वैतसृतिं विलुप्तविघ्नाम्}
{विदधत् प्रविशन् दरीं पुराणां मद् अघं हन्तु तिरोहितः स मौनी}% ॥ ६५ ॥

\twolineshloka
{सशरीरम् उपेयिवांसम् ओकः शशिचूडस्य जगद्गुरुं तम् एकम्}
{अपि चीनतुरुष्कबाह्लिकाद्यैः स्वपराचार्यतया स्तुतं स्तुमोऽद्य}% ॥ ६६ ॥

\twolineshloka
{अनुपद्मपुरं चिराय तिष्ठन् ध्वनिकारादिबुधाश्रितो व्रतिष्ठः}
{पुनरञ्चितकाञ्चिरस्तु पक्षे मम सच्चिच्छिखरो विलासभिक्षुः}% ॥ ६७ ॥

\twolineshloka
{अतिरूपम् अपापवृत्तमान्यं यतिताम् आश्रितम् अङ्गजं यथान्यम्}
{समयोच्चयसारसावधृत्यै सुमहादेवगुरुं स्मरामि धृत्यै}% ॥ ६८ ॥

\twolineshloka
{कृतसट्टकसत्रिनाट्यबन्धव्रतयायावरराजशेखरान्ध्यम्}
{हृतवन्तम् अनन्तमन्त्रशक्तिं व्रतिगङ्गाधरम् आश्रयेऽर्थ्यसूक्तिम्}% ॥ ६९ ॥

\twolineshloka
{शिवयोजितचन्द्रमौलिपूजाप्रवणं तन्नियमित्रयं महौजः}
{परवस्तुविद् अस्तसम्प्रचारादरम् आस्तां परम् अद्य मे विचारे}% ॥ ७० ॥

\twolineshloka
{कविसोमककॢप्तसेवनत्यग् भविदूरोऽद्रिगुहागृहेऽहिकृत्यः}
{चिरम् आत्मनि सञ्चरन् प्रमत्तः परमाद्यः स शिवः शिवं विधत्ताम्}% ॥ ७१ ॥

\twolineshloka
{नवभोजमहीपदत्तमुक्ताशिबिकाभ्रान्तसमस्तदक्षिणाध्वा}
{कलशेश्वरमन्त्रिणाऽस्तबाधं कलयन् स्वं मठम् ईड्यतां स बोधः}% ॥ ७२ ॥

\twolineshloka
{जयदेवसुमङ्खकृष्णमिश्राश्रयताम् एत्य कुमारपालमिश्रम्}
{अपि हैमम् अपाकृतोक्तियोगैरुपसेव्योऽस्तु च चन्द्रचूडयोगी}% ॥ ७३ ॥

\twolineshloka
{अभिचारकगुप्तपादवादिप्रभुहर्षादिपराभवाग्रभूमिम्}
{कलये हृदि संश्रितं स्वभासा विलयं चिद् वियतीह चिद्विलासम्}% ॥ ७४ ॥

\twolineshloka
{बहुधा कृतकोटिहोमलीलावभृथौ शाक्तपथाववामशीलौ}
{क्रमशः प्रणतोऽस्मि पावनाख्यौ सुमहादेवसुचन्द्रशेखराख्यौ}% ॥ ७५ ॥

\twolineshloka
{हृदि सायणमाधवोपदेष्टॄन् मदिनां वादविधौ शमप्रदेष्टॄन्}
{नमत श्रुतिशीलनाद् अतन्द्रान् यमिविद्याधिकतीर्थदेशिकेन्द्रान्}% ॥ ७६ ॥

\twolineshloka
{यमिषड्गवम् अप्लुतस्वदेशं शमिते म्लेच्छकुलैर्महत् प्रकाशे}
{अजहत् स्वमतं जहत् प्रचारं विजयायास्तु विवेकिनाम् उदारम्}% ॥ ७७ ॥

\twolineshloka
{समुपासितनीलकण्ठमन्त्रक्रमनिर्धूतविषामयं स्वतन्त्रम्}
{शिवयोगिसमाह्वयं यमीन्द्रं भवभेदाय भजाम मोदसान्द्रम्}% ॥ ७८ ॥

\twolineshloka
{लुलितागमिकच्छटोपदेशं ललितार्चाश्रितहस्तम् अस्तपाशम्}
{कलिताखिलवेदशास्त्रकोशं कलये प्रत्यगितं महःप्रकाशम्}% ॥ ७९ ॥

\twolineshloka
{खलवैष्णवतैन्त्रिणीककोलाहलहेलाहितगालिरुक्तिजालैः}
{शिवयन् वरदं चिरेण सौम्यः शिवकृन्मेऽस्तु स शङ्करः प्रणम्यः}% ॥ ८० ॥

\twolineshloka
{निजनीवृद् अवग्रहेतिखेदत्यजनेपालनृपालपूज्यपादः}
{स पुरो मम साधु सन्निधत्तां विपुलानन्दसदाशिवोऽप्रमत्तः}% ॥ ८१ ॥

\twolineshloka
{बहुधाऽऽहितसूत्रभाष्यशान्तिर्बहुलामोदयुतः समग्रशान्तिः}
{शमलं समुदस्य शं ममालं स महादेवगुरुर्दिशेत् सलीलम्}% ॥ ८२ ॥

\twolineshloka
{सतताहितचन्द्रमौलिसेवः श्रितकाञ्चीपुर एव शुद्धभावः}
{सृजतान्मम चन्द्रचूडमौनी स जयं सर्वत एव साधुमानी}% ॥ ८३ ॥

\twolineshloka
{अथ मे स सदाशिवः कृषीष्ट प्रथमं बोधपदेन तत्त्वदेष्टा}
{कुशलं कुशलाग्रणीः प्रकामप्रशमोऽप्युज्ज्वलमूर्तिराप्तकामः}% ॥ ८४ ॥

\twolineshloka
{हृदये स पदं सदा विधत्तां सदयं मे प्रणुदन् प्रमादवत्ताम्}
{निगमान्तगुरुः परः शिवात्मा सुगमः साधुततेर्यमी प्रशात्मा}% ॥ ८५ ॥

\twolineshloka
{यद् उदीक्षणलेशतो जनानां विदुरद्वैतमतिं द्वितां धुनानाम्}
{तम् अहं परमाञ्चितं शिवेन्द्रं शमवन्तं शरणं श्रये यमीन्द्रम्}% ॥ ८६ ॥

\twolineshloka
{इति नाथनयाऽऽत्मबोधनाम्नो यतिराजस्य सदाशिवेन्द्रभूम्ना}
{कलिता जयताद् वसन्तमालाषडशीत्या गुरुराजरत्नमाला}% ॥ ८७ ॥

{इति जगद्गुरुरत्नमाला सम्पूर्णा}
